\subsection{Excluding order $46$}\label{sec:d31a46}
\begin{proposition}\label{prop:decic46-exclusion}
An elliptic curve over a number field of degree $10$ has no point of order $46$.
\end{proposition}
\begin{proof}
Put $X=X_1(46)$ and $Y=X_1(23)$.
The curve $X$ has genus $45$ and rank-zero rational Jacobian
\cite[Theorem~3.1]{DEHMZ}.
Moreover, $\operatorname{gon}_{\mathbb Q}(X)=19$ by
\cite[Theorem~D]{Varivoda26}.
Let $D$ be a hypothetical noncuspidal effective rational divisor of degree $10$.

Entirely cuspidal reduction at $3$ lifts to an effective cuspidal divisor $C$
over $\mathbb Q_3$.
With a rational cusp $c_0$, both $[D-10c_0]$ and $[C-10c_0]$ are torsion.
Since the torsion kernel of reduction over $\mathbb Q_3$ is zero, $D\sim C$.
The divisors are distinct, so $h^0(X_{\mathbb Q_3},D)\ge2$.
Cohomology and base change give $h^0(X,D)\ge2$, contradicting rational gonality $19$.

The complete noncuspidal closed-point counts at $3$ are
\[
 (0,0,0,55,44,55,264,781,1958,6292).
\]
They are independently checked using Frobenius similarity classes.
The global and finite cusp signatures both equal $1^{22}11^2$.
In particular, every possible cuspidal complement of a nonempty
noncuspidal part is a sum of rational cusps.
For any good odd prime $q$ put
\[
 A_q=T_q-1-q\langle q\rangle.
\]
Eichler--Shimura and odd-prime torsion injectivity give $A_qJ_X(\mathbb Q)=0$.
The operator also annihilates a rational cusp as an actual divisor.
Thus $A_qD'$ must be principal for a noncuspidal part $D'$ whenever it is used below.

We also use the degree-$3$ forgetful map
$\pi:X\longrightarrow Y$, which doubles the point of order $46$.
By \cite[Theorems~3.1 and~4.13, Table~2]{DEHMZ},
$J_Y(\mathbb Q)$ is cuspidally generated of order $408991$.
The proven reduction group at $3$ is cyclic of order $9406793$.
Its quotient by the rational image is therefore
\[
 J_Y(\mathbb F_3)/\operatorname{red}_3J_Y(\mathbb Q)\simeq C_{23}.
\]
Every target cusp of degree at most $10$ is rational, so all allowed
cuspidal complements have zero image in this quotient.
Consequently the normalized pushforward class of $D'$ must vanish.
Residue-field degrees are included in each pushforward.

The retained tests in degrees at most $9$ give $782$ negative results,
with support types
\[
 (4),(5),(6),(7),(8),(4,4),(9),(5,4).
\]
They remain valid here: the Hecke divisors are unchanged, and the larger
cuspidal complements still lift globally and vanish in the target quotient.
The only new support types are $(10)$, $(6,4)$ and $(5,5)$.
The exact test counts and quotient results are
\[
\begin{array}{c|r|r|r}
\text{type}&\text{tests}&\text{quotient rejections}&\text{Hecke tests}\\\hline
(10)&572&552&20\\
(6,4)&275&264&11\\
(5,5)&176&170&6
\end{array}
\]
For the mixed types we normalize the first factor by diamonds and retain
every second factor, including repeated places.

The quotient computation is replayed independently on all $1023$ supports.
It uses $676$ marked source components.
Their elliptic equations, exact level orders, actual doubled points,
residue degrees and norms are checked separately.
An independent finite-field reconstruction of the diamond orbits gives
all $6292$ decic source places and the required smaller support sets.
The complete class sums give exactly $986$ nonzero images and $37$ exceptions.

For each exception, an exact computation on the genus-$45$ source gives
\[
 h^0\bigl(X_{\mathbb F_3},A_5D'\bigr)=0.
\]
The complete cyclic $5$-kernel partition is retained, with the actual
point of order $46$ transported through every isogeny.
A separate native replay reconstructs the prime ideals and degree-zero
Hecke divisors and again gives dimension zero in all $37$ cases.
Independent V\'elu computations verify all $324$ geometric kernels of
the $54$ marked elliptic components, including every residue-field norm.
Thus none of the exceptional divisors is principal.
All possible reductions of $D$ have been excluded.
\end{proof}
